% ============================================================================
\clearpage
\section*{v4.9.0-r72 Source-Only: The \texorpdfstring{\(\vartheta\)}{theta} Gauge of the OP5 Comparison Symbol, a Cartan Test, and Four Hidden Conditions}
\addcontentsline{toc}{section}{v4.9.0-r72 Source-Only: The theta Gauge of the OP5 Comparison Symbol, a Cartan Test, and Four Hidden Conditions}
\label{sec:v490-r72-theta-gauge}
% ============================================================================

\paragraph{Purpose.}
Layer v4.9.0-r71 found that the free function of the \(\Omega\)-augmented OP5 comparison symbol is an exact gauge symmetry
\(\delta\Psi=f\,\partial_\vartheta\Psi\) and named the next step: fix the \(\vartheta\) gauge, then test involution.  On
3 November 2026 the author asked for exactly this.  This layer records the gauge, one confirmed pre-registered Cartan test
(H51), three further pre-registered tests that turned out invalid by their own criteria (H52, H53, H54), and one exploration of
the Spencer cohomology of the gauged symbol.  Everything concerns the comparison model of \S\ref{sec:v410-state-mediator}
only.  Earlier layers are not rewritten.

\noindent\textbf{Source-only status.}
This v4.9.0-r72 layer is source-only.  It does not build a PDF, does not assign a new DOI,
and does not authorize a public release.  The v4.9.0 Version DOI
\href{https://doi.org/10.5281/zenodo.23534954}{10.5281/zenodo.23534954} and the Concept DOI
\href{https://doi.org/10.5281/zenodo.19210728}{10.5281/zenodo.19210728} are unchanged by it.

\subsection*{v4.9.0-r72 results}
\label{subsec:v490-r72-results}

Notation as in \S\ref{sec:v490-r71-op5-symbol}: \(R=(F,G,dg)\), \(J\) the completed first-order symbol, \(\theta\) the
symbol of \(\Omega\), \(g_p\) the degree-\(p\) symbol of the \(\Omega\)-augmented system.  \(g'_p\) is the symbol after the
gauge condition of r72.1, \(c_k\) denotes the Taylor coefficients of degree \(k\) of \(\Psi\), and \(P_{5}\), \(P_{56}\),
\(P_{567}\) are the dimensions of the projections of a solution space of the linearized jet tower onto \(c_5\),
\((c_5,c_6)\), \((c_5,c_6,c_7)\).

\paragraph{Theoremlet v4.9.0-r72.1 (the \(\vartheta\) gauge).}
\begin{enumerate}[leftmargin=2em]
\item[(a)] The gauge vector \(v=\partial_\vartheta\Psi\) at the germ is a common kernel vector of all four blocks of \(J\)
(smallest singular value \(5.5\times10^{-15}\) against \(0.13\)), and all blocks of \(\theta\) annihilate it
(\(1.6\times10^{-14}\)).
\item[(b)] The algebraic condition \(s(\Psi)=S(x)\) with \(s=\Psi^\dagger K_2\Psi\), one component of the current, is
transversal: \(|ds\cdot v|/|ds|=0.45\), the largest value among \(\Psi^\dagger\Psi\), \(\Psi^\dagger K_a\Psi\),
\(\Psi^\dagger Q_a\Psi\), \(\Psi^\dagger(I\otimes\beta)\Psi\).  \(\Psi^\dagger(Z\otimes\beta)\Psi\) is fixed by \(g\) and blind
(\(3.7\times10^{-16}\)).  Because \(ds\cdot v\neq0\), the condition fixes \(f\) locally.
\item[(c)] At symbol level every transversal algebraic gauge gives \(g'_p\cong g_p/(S^p\otimes v)\).  Measured with clean
gaps: \(g'_p=g_p-\binom{p+3}{3}=117,188,273,371,481,606,745\) for \(p=2,\dots,8\).
\end{enumerate}

\paragraph{Theoremlet v4.9.0-r72.2 (H51 confirmed: the gauged symbol is involutive at order 6).}
Pre-registered before the computation.  Cartan's test \(\dim g'_{p+1}\le\sum_{k=0}^{3}\dim g'^{(k)}_p\) in three random
orthonormal frames, all ranks with gap ratio at least \(10^7\):
\begin{center}
\small
\begin{tabular}{@{}ccccc@{}}
\toprule
\(p\) & \(\dim g'^{(k)}_p\), \(k=0..3\) & sum & \(\dim g'_{p+1}\) & verdict\\
\midrule
4 & 273, 85, 14, 0 & 372 & 371 & not involutive\\
5 & 371, 98, 14, 0 & 483 & 481 & not involutive\\
6 & 481, 111, 14, 0 & 606 & 606 & involutive\\
7 & 606, 125, 14, 0 & 745 & 745 & involutive\\
\bottomrule
\end{tabular}
\end{center}
After the gauge no free function of four variables remains; formally the symbol leaves 14 real functions of three variables.

\paragraph{Record v4.9.0-r72.3 (H52 invalid; the stored fifth jet does not extend).}
Pre-registered test of formal integrability from order 6 with the stored jet through degree 5 held fixed.  The ranks in the
large systems had no clean gap (ratio \(70\) to \(600\)); the test is invalid by its criterion.  A clean subsystem (gap
\(6\times10^7\)) shows that the stored degree-5 coefficients cannot be extended at all: the compatibility residual is
\(4.0\times10^{-3}\), at the level of the positive control.  They were solved on 30 October only against the degree-4
equations and violate the once prolonged \(\Omega\) condition.  The counts agree: without the prolonged \(\Omega\) rows the
fibres have \(381=371+10\) and \(501=481+20\) dimensions.  This was an error of the test design, not a property of the system.

\paragraph{Record v4.9.0-r72.4 (H53 invalid; the tower over the repaired fourth jet).}
Corrected design: only the repaired jet through degree~4 is fixed, \(c_5,\dots,c_N\) are free, and the linearization
\(L=A^\mu(x)\partial_\mu+B(x)\) is exact (symmetric probes \(x^{m e_\nu}e_i\); even terms cancel, cubic terms start above the
degrees used).  Reference residual \(4.6\times10^{-14}\), positive control \(3.5\times10^{-3}\).
\begin{enumerate}[leftmargin=2em]
\item[(a)] With equations through degree 6 (\(F_7\)) everything is clean: residual \(4.0\times10^{-12}\), \(P_5=371\),
\(P_{56}-P_5=481\).  Up to these equations there are no conditions beyond \((F,G,dg)\), \(\Omega\) and the gauge.
\item[(b)] With equations through degree 7 (\(F_8\), \(9865\times10200\)) \(P_{56}\) drops from \(852\) to \(848\), while
\(P_{567}-P_{56}=606=\dim g'_7\).  One rank had gap ratio \(4\times10^3\); the test is invalid by its criterion.
\end{enumerate}

\paragraph{Theoremlet v4.9.0-r72.5 (exploration: Spencer cohomology of the gauged symbol).}
Not pre-registered.  With \(H^{k,2}=\dim\ker(\delta:g_k\otimes\Lambda^2\to S^{k-1}\otimes\Lambda^3)-\operatorname{rank}(\delta:g_{k+1}\otimes\Lambda^1\to S^k\otimes\Lambda^2)\),
all gaps at least \(4\times10^6\): \(H^{4,2}=0\) and \(H^{5,2}=4\) for the ungauged symbol (check against r71.2) and
\(H^{4,2}(g')=0\), \(H^{5,2}(g')=4\) for the gauged one.  At most four hidden conditions on 6-jets are possible, none on 5-jets.
The gauge does not change this, so the four do not come from the \(\vartheta\) freedom.

\paragraph{Record v4.9.0-r72.6 (H54 invalid; an error bound says exactly four).}
Pre-registered.  By Weyl's inequality an error \(\Delta E\) moves every singular value by at most \(\|\Delta E\|_2\).  Measured
from independent estimates: the reduced matrix behind \(P_{56}\) contains no degree-3 coefficients of \(A,B\) and has error
\(\le7\times10^{-13}\) against its smallest counted value \(1.0\times10^{-9}\); the full \(F_8\) has error about
\(5\times10^{-9}\) against its smallest nonzero value \(2.7\times10^{-6}\), and two perturbed variants keep rank \(7957\).
Hence \(P_{56}\le848\), and with r72.5 \(P_{56}\ge848\).  The criterion also demanded a gap ratio of \(10^4\) in the perturbed
variants, which the perturbation itself makes impossible (ratios \(3\times10^3\), \(7\times10^3\)); the test is invalid.  The
pre-registered variant (degree-3 part of \(B\)) does not enter \(F_8\) at all; a stricter substitute was used and recorded.

\paragraph{Proof and measurement sources.}
r72.1 by \texttt{k136\_eichung.py}.  r72.2 by \texttt{k137\_cartan\_test.py}.  r72.3 by \texttt{k138\_tps\_d.py} and
\texttt{k139\_fortsetzung.py}.  r72.4 by \texttt{k140\_fortsetzung\_exakt.py} and \texttt{k140\_f8\_schritte.py}.  r72.5 by
\texttt{k141\_spencer\_geeicht.py}.  r72.6 by \texttt{k142\_h54.py}.  The notes \texttt{Ergebnis\_H51\_\allowbreak
Cartan\_\allowbreak Test\_\allowbreak geeichtes\_\allowbreak Symbol.md} to \texttt{Ergebnis\_H54\_\allowbreak vier\_\allowbreak
Bedingungen\_\allowbreak 6Jets.md} and \texttt{Erkundung\_\allowbreak Spencer\_\allowbreak geeichtes\_\allowbreak Symbol.md}
give the numbers.  \(\square\)

\paragraph{Consequence.}
The gauged symbol is involutive from order 6 (H51).  The gauged system itself is, by all evidence, not formally integrable at
order 6: four further conditions on 6-jets appear once the degree-7 equations are used (r72.4(b), r72.5, r72.6), but no test
has confirmed this formally.  A confirming test needs new data, for example a second base point.  If confirmed, the four
conditions must be determined, added to the system, and the Cartan test repeated.  OP5 is not closed.

\subsection*{v4.9.0-r72 status in the three frames}
\label{subsec:v490-r72-frames}

\begin{itemize}[leftmargin=2em]
\item All statements concern the fixed OP5 comparison model of \S\ref{sec:v410-state-mediator} in one chart with one coframe
reconstructed from \(\Psi\).  No frame map M, T or P acts on them; r72.1 to r72.6 are frame-free within the model.
\item The gauge function \(\Psi^\dagger K_2\Psi\) is a current component in the model coframe, not a frame-resolved quantity.
\end{itemize}
No statement of this layer concerns colour or charge.

\subsection*{v4.9.0-r72 register entries}
\label{subsec:v490-r72-register}

\begin{center}
\small
\begin{tabular}{@{}>{\raggedright\arraybackslash}p{0.34\textwidth}>{\raggedright\arraybackslash}p{0.58\textwidth}@{}}
\toprule
statement & status after v4.9.0-r72\\
\midrule
\(\vartheta\) gauge & algebraic, \(\Psi^\dagger K_2\Psi=S(x)\); removes exactly the free function\\
H51 & confirmed: gauged symbol involutive first at \(p=6\)\\
H52, H53, H54 & invalid by their criteria (test design); findings recorded\\
stored degree-5 coefficients & not extendable; degree 5 must stay free\\
\(H^{5,2}(g')\) & \(4\) (exploration)\\
hidden conditions on 6-jets & by all evidence exactly four; not formally confirmed\\
next step & second base point (H55), then the explicit conditions and a new Cartan test\\
\bottomrule
\end{tabular}
\end{center}

\subsection*{v4.9.0-r72 ledger and guard update}
\label{subsec:v490-r72-ledger-guard}

\[
R_{\rm chk}^{v490\text{-}r72}=(1,1,0)
\]
for (gauge fixed and H51 confirmed; H52 to H54 recorded as invalid; OP5 not closed).

The layer spends no guard:
\[
R_{\rm glob}^{(v490\text{-}r72,+)}=(0,1,1,1,1),
\qquad
\Gamma_{\rm guard}^{v490\text{-}r72}=(1,1,1,1,1),
\qquad
\beta_4=1/256\ \text{unspent}.
\]
It does not imply exact finite \(C5\) zero closure, exact all-readout terminality,
\(C6/T2\), a global photon-clock/RH theorem, or global all-readout closure.
